%%
%% paper22_dag.tex  --  Logical Dependency Graph for Paper XXII
%%
%% Every edge A -> B: ``B requires A in its proof''.
%% Format is compatible with trilogy_dag.tex (XVIII--XX).
%%
%% Structure:
%%   Layer 0: Imported black boxes (from other papers in the program)
%%   Layer 1: IIKS representation (§2)
%%   Layer 2: Two-level obstruction (§3)
%%   Layer 3: Gram-to-gap reduction (§4)
%%   Layer 4: Airy closure (§5)
%%   Layer 5: Open problems with mechanisms
%%
\documentclass[12pt,a4paper]{amsart}
\usepackage{amsmath,amssymb,amsthm,mathtools,enumitem,booktabs,hyperref}

\newtheorem{theorem}{Theorem}[section]
\newtheorem{lemma}[theorem]{Lemma}
\newtheorem{proposition}[theorem]{Proposition}
\theoremstyle{definition}
\newtheorem{definition}[theorem]{Definition}
\newtheorem{remark}[theorem]{Remark}

\begin{document}

\title{Logical DAG: Paper~XXII (Paper-Grade)\\
Edge Coercivity and the Airy Boundary Operator}
\date{May 2026}
\maketitle

\tableofcontents

%%=================================================================
\section{Purpose and Conventions}
\label{sec:purpose}
%%=================================================================

This document is the paper-grade logical dependency graph for Paper~XXII.
It follows the format of \texttt{trilogy\_dag.tex} (XVIII--XX).
Every edge is explicit; every open edge carries a constructive mechanism.

\medskip
\textbf{Conventions.}
(T) = proved theorem within XXII or cited paper.
(I) = exact identity (no asymptotics).
(H) = hypothesis / assumption (open).
(C) = conjecture (stated, not proved).
(BB) = black box (imported from another paper, no re-proof in XXII).
(O) = open problem with explicit mechanism.
An edge $A\to B$ means: B's proof invokes A.
An edge $A\xrightarrow{\mathrm{mech}}B$ additionally names
the explicit transfer mechanism.

\medskip
\textbf{Scope of XXII.}
The paper has exactly one new open question:
\[
  \inf\,\sigma(\Dedge) > 0
  \quad (\text{XXII Main Conjecture})
\]
where $\Dedge = \Pi_{\mathrm{edge}}(A_{N,c}-\mathcal{K}_c)\Pi_{\mathrm{edge}}$.
All other statements are either proved theorems, imported black boxes,
or explicitly labelled conditionals.

%%=================================================================
\section{Layer~0: Imported Black Boxes}
\label{sec:L0}
%%=================================================================

These results are used in XXII but not re-proved.
Every BB-edge is a theorem in the cited paper.

\begin{center}
\renewcommand{\arraystretch}{1.35}
\begin{tabular}{l p{5.5cm} l l}
\toprule
\textbf{ID} & \textbf{Statement} & \textbf{Source} & \textbf{Used in} \\
\midrule
BB1 & $\sigma(A_{N,c})\approx\sigma(\mathcal{K}_c)$ in bulk
    (Theorem~A)
    & Paper~XXI
    & P1 \\
BB2 & $\Dedge\ne 0$: at least one positive eigenvalue
    (Theorem~B)
    & Paper~XXI
    & P1, P5 \\
BB3 & Gram coercivity $\Rightarrow$ spectral gap
    $\ge Cc^{-1/2}$ (Theorem~A)
    & Paper~13b
    & P9 \\
BB4 & $K_c(x,y)=\frac{\sin(c(x-y))}{\pi(x-y)}$:
    sinc kernel representation
    & Slepian 1978
    & P2 \\
BB5 & Widom edge phase $\phi_c$
    and $O(c)$ growth
    & Widom 1964
    & P6, O3 \\
BB6 & Deift--Zhou steepest descent framework
    & Deift 1999
    & P7, P8 \\
BB7 & Fredholm determinant continuity
    in $\mathfrak{S}_1$-norm
    & standard
    & P10 \\
\bottomrule
\end{tabular}
\end{center}

%%=================================================================
\section{Layer~1: IIKS Representation (\S2)}
\label{sec:L1}
%%=================================================================

\begin{center}
\renewcommand{\arraystretch}{1.35}
\begin{tabular}{l p{5.5cm} l l}
\toprule
\textbf{ID} & \textbf{Statement} & \textbf{Type} & \textbf{Used in} \\
\midrule
P1 & Operator setup: $\Dedge$
    as mismatch form, Bulk vs.\ Edge
    & (T), from BB1, BB2
    & P3, P4 \\
P2 & Exact IIKS identity:
    $K_c(x,y) = \vec{F}_c(x)^\top J\,\vec{G}_c(y)/(x-y)$
    & (I), from BB4
    & P3, P7 \\
P3 & The cubic Airy term $\tfrac{1}{3}u^3$
    does not arise from the Fourier
    representation; it requires
    integration of the Widom edge phase
    via the IIKS RHP structure
    & (T) structural remark, from P2, BB5
    & P6, P7 \\
\bottomrule
\end{tabular}
\end{center}

\begin{remark}[Why P2 is infrastructure, not argument]
P2 is an exact algebraic identity.
Its role in XXII is not to prove anything about $\Dedge$
directly, but to provide the coordinate system in which
Assumptions~A and~B (Layer~2) can be formulated.
Without P2, the RHP associated to $\mathcal{K}_c$ has no
well-defined jump matrix, and the obstruction decomposition
(P5--P7) cannot be stated.
\end{remark}

%%=================================================================
\section{Layer~2: Two-Level Obstruction Principle (\S3)}
\label{sec:L2}
%%=================================================================

\begin{center}
\renewcommand{\arraystretch}{1.35}
\begin{tabular}{l p{5.5cm} l l}
\toprule
\textbf{ID} & \textbf{Statement} & \textbf{Type} & \textbf{Used in} \\
\midrule
P4 & Edge stability decomposes as:
    (nonlinear phase consistency) $+$
    (linear kernel transfer)
    & (T) structural, from P1, P3
    & P5, P6, P7 \\
P5 & Assumption~A: $\exists\, g_c$ analytic
    satisfying nonlinear matching
    $2\mathrm{Re}(g_c^+) = \mathrm{Re}(\phi_c)$
    with $g_c(z)=\log z + O(z^{-1})$
    & (H) open, from BB5
    & P7, P8 \\
P6 & Nonlinearity of Assumption~A:
    $\phi_c = \phi_c[g_c]$ (implicit
    dependence through spectral density
    near band edge); no RMT-type
    candidate satisfies (A) for large $c$
    & (T) structural, from P3, BB5
    & P5, O3 \\
P7 & Assumption~B: given Assumption~A,
    steepest-descent geometry holds
    (sign condition, Lipschitz contours,
    no additional saddles)
    & (H) conditional on P5
    & P8 \\
P8 & Conjecture~3.3 (RHP convergence):
    $Y_c\to Y_{\mathrm{Ai}}$ uniformly
    on compacta; implies
    $\|\tilde{K}_c^{(\xi)}-K^{\mathrm{Airy}}\|_{\mathfrak{S}_1}=O(c^{-1/3})$
    & (C) conditional on P5, P7, BB6
    & P10, P11 \\
\bottomrule
\end{tabular}
\end{center}

\begin{remark}[Logical isolation of A and B]
P5 (Assumption~A) is a nonlinear existence statement;
it does not follow from P7 and cannot be reduced to it.
P7 (Assumption~B) is a geometric regularity statement
that is standard in form but cannot even be formulated
without P5, since $2g_c - \phi_c$ is undefined without
a compatible $g_c$.
Thus the logical direction is always $P5 \to P7$, never
$P7 \to P5$.
The deep difficulty lies entirely in P5.
\end{remark}

%%=================================================================
\section{Layer~3: Gram-to-Gap Reduction (\S4)}
\label{sec:L3}
%%=================================================================

\begin{center}
\renewcommand{\arraystretch}{1.35}
\begin{tabular}{l p{5.5cm} l l}
\toprule
\textbf{ID} & \textbf{Statement} & \textbf{Type} & \textbf{Used in} \\
\midrule
P9 & Reduction: uniform edge-frame
    stability $\Rightarrow$
    $\lambda_{\min}(G_{N,c})\ge\alpha(c)
    \asymp c^{-1/2}$
    (Prop.~4.1)
    & (T), from Landau--Widom
    & P12 \\
R1 & (Reduction functor via BB3):
    $\lambda_{\min}(G_{N,c})\ge c^{-1/2}$
    $\Rightarrow$ gap lower bound
    $\ge Cc^{-1/2}$ (Theorem~4.2)
    & (BB) = BB3 imported, from P9
    & P12 \\
P10 & Conditional Airy-scale upgrade:
    under Conj.~3.3, gap upgrades
    to $c^{-1/3}$ (Corollary~4.1)
    & (C) conditional on P8, BB7
    & P13 \\
\bottomrule
\end{tabular}
\end{center}

\begin{remark}[R1 is not re-proved]
R1 imports BB3 (Paper~13b Theorem~A) exactly.
The new content of XXII is P9 (the first arrow in the chain):
showing that uniform quasi-isometry of $\mathrm{Ev}|_{\mathrm{Edge}}$
controls the Gram minimum eigenvalue at the required scale.
R1 is then applied as a black-box functor.
XXII does not touch the proof of BB3.
\end{remark}

%%=================================================================
\section{Layer~4: Airy Closure and Main Theorem (\S5)}
\label{sec:L4}
%%=================================================================

\begin{center}
\renewcommand{\arraystretch}{1.35}
\begin{tabular}{l p{5.5cm} l l}
\toprule
\textbf{ID} & \textbf{Statement} & \textbf{Type} & \textbf{Used in} \\
\midrule
P11 & $D^{\mathrm{Airy}} =
     \Pi_{\mathrm{Airy}}(A^{\mathrm{arith}}
     - K^{\mathrm{Airy}})\Pi_{\mathrm{Airy}}$:
     Rand-Weyl operator definition
     & (T) from P8, Prop.~5.1
     & P12 \\
P12 & Conjecture~5.2 (main gap conjecture):
     $\inf\,\sigma(D^{\mathrm{Airy}})
     \ge c_2^{\mathrm{Airy}} > 0$
     & (C) open
    & P13 \\
P13 & Main Theorem (Thm.~5.3,
     conditional on P12):
     $\inf_{f\in\mathrm{Edge},\|f\|=1}
     \langle\Dedge f,f\rangle
     \ge c_2^{\mathrm{Airy}}/2 > 0$
     & (T) conditional on P11, P12,
       from P9, R1, P8
     & (terminal) \\
\bottomrule
\end{tabular}
\end{center}

\begin{remark}[Airy closure is asymptotic, not foundational]
P11--P13 are not the logical ground of XXII.
They are the asymptotic ceiling: the scaling limit into which
the reduction chain (P9, R1) feeds.
The Airy kernel does not justify the coercivity;
it parametrizes the limit operator whose spectral gap
is the one remaining unknown.
\end{remark}

%%=================================================================
\section{Layer~5: Open Problems with Mechanisms}
\label{sec:L5}
%%=================================================================

\begin{center}
\renewcommand{\arraystretch}{1.35}
\begin{tabular}{l p{3.8cm} l p{4.5cm}}
\toprule
\textbf{ID} & \textbf{Statement} & \textbf{Closes} & \textbf{Mechanism} \\
\midrule
O1 & Prove Assumption~A:
    $\exists\,g_c$ satisfying
    \eqref{eq:g-matching}
    & P5
    & Fixed-point operator
      $\mathcal{T}_c(g)
      = \mathcal{H}_{(-1,1)}
      [\tfrac{1}{2}\mathrm{Re}(\phi_c[g])]
      + \log(\cdot)$;
      prove $\mathcal{T}_c$ contractive
      on a suitable Banach space
      of analytic functions;
      standard: Muskhelishvili theory
      once nonlinear coupling is
      controlled \\
O2 & Prove Assumption~B
    given~A
    & P7
    & Standard Deift--Zhou:
      sign check on $\mathrm{Re}(2g_c-\phi_c)$
      along contour, Lipschitz
      parameterisation;
      no new method required
      once $g_c$ is known from O1 \\
O3 & Uniform BW-Doubling:
    $\inf_{n\in[(1-\delta)N,N)}
    \mathcal{E}_{\mathrm{out}}(|\psi_n|^2)
    \ge c_3 > 0$
    & O4 (alternate route to P9)
    & Extend Paper~III pointwise
      argument to a uniform
      Schur-test on the edge block;
      alternatively: transfer
      via Kadec-$\tfrac{1}{4}$ for
      prime frequencies in edge band \\
O4 & Uniform quasi-isometry:
    $A\|f\|^2\le\|\mathrm{Ev}\,f\|^2
    \le B\|f\|^2$ for all
    $f\in\mathrm{Edge}$
    & P9 (direct route)
    & Frame stability via
      prime-frequency sampling
      in edge band;
      equivalent to Kadec-type
      stability for $\mathrm{PW}_c$
      restricted to edge \\
O5 & Compute $c_2^{\mathrm{Airy}}$
    explicitly
    & P12
    & Quantitative oscillation
      bounds on $\mathrm{Ai}(t)$;
      $\mathrm{Ai}(t)>0$ for $t>0$
      implies non-degenerate
      prime sampling of
      $|\mathrm{Ai}(t)|^2$;
      connection to $a_1\approx-2.338$
      via spectral counting \\
O6 & Under RH: characterise
    RH-collapse of the
    coercivity gap
    & P13 (long-range)
    & $\sigma(D^{\mathrm{Airy}})$
      controlled by zero distribution
      of $\zeta(s)$; coercivity
      failure $\Leftrightarrow$
      accumulation of zeros off
      critical line \\
\bottomrule
\end{tabular}
\end{center}

%%=================================================================
\section{Critical Path Analysis}
\label{sec:critical_path}
%%=================================================================

\subsection{Minimal path to the unconditional gap bound}
\[
  \text{O4}
  \xrightarrow{\text{frame stability}}
  \text{P9}
  \xrightarrow{\text{Landau--Widom}}
  \lambda_{\min}(G_{N,c})\ge c^{-1/2}
  \xrightarrow{\text{BB3 (Paper~13b)}}
  \text{gap}\ge Cc^{-1/2}.
\]
This path is fully explicit once O4 is resolved.
O4 is independent of the RHP machinery (Layers 1--2).

\subsection{Minimal path to the Airy-scale upgrade ($c^{-1/3}$)}
\[
  \text{O1}
  \xrightarrow{\text{Muskhelishvili}}
  \text{P5}\,(\mathrm{Ass.~A})
  \xrightarrow{}
  \text{O2}
  \xrightarrow{\text{Deift--Zhou}}
  \text{P7}\,(\mathrm{Ass.~B})
  \xrightarrow{\text{BB6}}
  \text{P8}\,(\mathrm{Conj.~3.3})
  \xrightarrow{\text{BB7}}
  \text{P10}\,(c^{-1/3}\text{ gap}).
\]
This path is the deep route and requires the resolution of
the nonlinear RHP existence problem O1.

\subsection{Minimal path to $\inf\,\sigma(\Dedge)>0$ (Main Theorem)}
\[
  \text{O4} + \text{O5}
  \;\Rightarrow\;
  \text{P9} + \text{P12}
  \;\Rightarrow\;
  \text{R1} + \text{P13}\;\;(\mathrm{Main~Theorem}).
\]
O4 and O5 are independent.
O5 does not require O1 or O2.

\subsection{Critical observation}

The XXII architecture has \emph{two independent open inputs},
not one:
\begin{enumerate}[label=(\roman*),nosep]
  \item \textbf{O4} (uniform quasi-isometry at the edge):
    purely functional-analytic, no RHP involved.
  \item \textbf{O5} (spectral gap of $D^{\mathrm{Airy}}$):
    analytic, Airy function oscillation, no arithmetic involved.
\end{enumerate}
They feed the Main Theorem via disjoint sub-chains.
Resolution of either one in isolation yields a partial result;
both are required for the full theorem.
The RHP route (O1, O2) provides the deeper structural picture
(Airy-scale gap, $c^{-1/3}$ exponent) but is not on the
minimal path to $\inf\,\sigma(\Dedge)>0$.

%%=================================================================
\section{Dependency Graph (Textual)}
\label{sec:dag_textual}
%%=================================================================

\begin{verbatim}
[BB4: Slepian sinc kernel]
         |
         v
[P2: exact IIKS identity] ----+
         |                    |
         v                    |
[P3: cubic term = Widom      |
     phase integral]         |
         |                    |
         v                    v
[BB5: Widom edge phase]   [BB6: Deift-Zhou]
         |                    |
         v                    |
[P6: nonlinearity             |
     of phi_c]                |
         |                    |
         v                    |
[P5: Assumption A]            |
    (OPEN: O1) <-----------+  |
         |                 |  |
         v                 |  |
[P7: Assumption B]         |  |
    (OPEN: O2) <---------+ |  |
         |               | |  |
         +---------+-----+ |  |
                   |       |  |
                   v       |  |
[BB6+BB7]---->[P8: Conj.3.3: Y_c -> Y_Ai]
                   |
                   v
[P10: c^{-1/3} gap scale]

----------- Independent sub-chain -----------

[BB1+BB2: Paper XXI]
         |
         v
[P1: Dedge setup]
         |
         v
[P4: Edge stability decomposition]
         |
         +-----------+
         |           |
         v           v
  [P9: Gram]    [P11: D^Airy]
  (OPEN: O4)   (from P8)
         |           |
         v           v
  [R1 = BB3: Gram->Gap]
         |    +-------+
         |    |       |
         v    v       v
[P12: inf sigma(D^Airy)>0]  [P10]
         (OPEN: O5)
         |
         v
[P13: Main Theorem]
    (conditional on P12 + P9 + R1)
\end{verbatim}

%%=================================================================
\section{Freeze Checklist}
\label{sec:freeze}
%%=================================================================

\begin{center}
\renewcommand{\arraystretch}{1.3}
\begin{tabular}{l l}
\toprule
\textbf{Item} & \textbf{Status} \\
\midrule
All BB0--BB7 sourced and not re-proved & \checkmark \\
P2 (IIKS) verified as exact identity & \checkmark \\
P3 (cubic term from Widom phase) stated with mechanism & \checkmark \\
P4 (edge stability decomposition) structurally derived & \checkmark \\
P5 (Assumption~A) explicitly labelled (H), not proved & \checkmark \\
P6 (nonlinearity of $\phi_c$) explicitly stated & \checkmark \\
P7 (Assumption~B) conditional on P5 only & \checkmark \\
P8 (RHP conjecture) labelled (C) with precise statement & \checkmark \\
P9 (Gram lower bound) with Landau--Widom scale & \checkmark \\
R1 (BB3) imported as functor, not re-proved & \checkmark \\
P10 (conditional $c^{-1/3}$ upgrade) labelled (C) & \checkmark \\
P11 (Rand-Weyl operator) defined from P8 & \checkmark \\
P12 (main gap conjecture) labelled (C) & \checkmark \\
P13 (Main Theorem) conditional scope in statement & \checkmark \\
All O1--O6 carry explicit mechanisms & \checkmark \\
Two independent open inputs (O4, O5) identified & \checkmark \\
RHP path (O1, O2) not on minimal path to Main Thm & \checkmark \\
\hline
O1: Assumption~A existence (nonlinear RHP) & open \\
O2: Assumption~B geometry & open (conditional on O1) \\
O3: Uniform BW-Doubling & open (alternate route) \\
O4: Uniform quasi-isometry $\Ev|_{\mathrm{Edge}}$ & open \\
O5: $c_2^{\mathrm{Airy}} = \inf\,\sigma(D^{\mathrm{Airy}}) > 0$ & open \\
O6: RH collapse characterisation & open (long-range) \\
\bottomrule
\end{tabular}
\end{center}

\medskip
\textbf{Freeze verdict.}
Paper~XXII is structurally complete as a conditional programme:
all logical dependencies are explicit, all open inputs are
identified with constructive mechanisms, and no BB is re-proved.
The minimal path to the Main Theorem requires O4 and O5,
which are independent and carry different methods
(functional-analytic vs.\ Airy oscillation).
The RHP route (O1--O2) deepens the result to the
$c^{-1/3}$ exponent but is not a prerequisite.

\begin{thebibliography}{99}
\bibitem{PaperXXI}
S.~Schmalnauer, \textit{Commutativity Failure in the
Arithmetic--Spectral Diagram} (Paper~XXI), preprint, 2026.
\bibitem{Paper13b}
S.~Schmalnauer, \textit{Spectral Gap Lower Bounds via
Gram Coercivity} (paper13b), preprint, 2026.
\bibitem{Paper13cc}
S.~Schmalnauer, \textit{Airy Universality for PSWF:
Exact IIKS Structure and RHP Reduction} (paper13cc), preprint, 2026.
\bibitem{Widom1964}
H.~Widom, \textit{Asymptotic behavior of the eigenvalues~II},
Arch.\ Rational Mech.\ Anal.\ \textbf{17} (1964).
\bibitem{Deift1999}
P.~Deift, \textit{Orthogonal Polynomials and Random Matrices:
A Riemann--Hilbert Approach}, AMS, 1999.
\bibitem{Muskhelishvili1953}
N.~Muskhelishvili, \textit{Singular Integral Equations},
Noordhoff, 1953.
\end{thebibliography}

\end{document}
